% day1_seg4_prob/moments.tex
\begin{frame}
  \onslide<1->{
	\frametitle{Statistics} 
	\framesubtitle{Moments: Means and Variances of Random Variables}
}
\begin{itemize}
  \onslide<2->{
\item 1st Order Moment: Expectation or Mean  
}
  \onslide<3->{
\item  2nd Order Moment: Variance 
}
  \onslide<4->{
\item Higher Order Moments  
}
\end{itemize}
\end{frame}
\begin{frame}
  \onslide<1->{
	\frametitle{Statistics} 
	\framesubtitle{1st Order Moment: Expectation of a Random Variable}
}
\begin{itemize}
  \onslide<2->{
 \item Let $X$ be a simple random variable and let
 $a_1,\ldots,a_n$ be the distinct values taken on by $X$.  Then,
 the expectation operator $\text{E}:SRV\rightarrow \Real$ is defined as
\begin{equation}\label{simp_expect}
  \expect{X}=\sum_{i=1}^{n}a_iP(A_i)
\end{equation}
where $A_i=X^{-1}\set{a_i}=\set{\omega\in\Omega:X(\omega)=a_i}$.
}
  \onslide<3->{
 \item Thus, the expectation of $X$ is the weighted
 sum of the values of $X$, the weights being the "size" or likelihood
 of the events.
}
  \onslide<4->{
 \item It is clear that $X=Y \text{a.s.}\Rightarrow
 \expect{X}=\expect{Y}$.
}
\end{itemize}
\end{frame}
\begin{frame}
  \onslide<1->{
	\frametitle{Statistics} 
	\framesubtitle{Expectation}
}
\begin{itemize}
  \onslide<2->{
 \item Alternative Integral notation:
$\expect{X}=\int_{\Omega}XdP$, because the weighted sum given in
\ref{simp_expect} is just the definition of the integral of $X$
with respect to the (probability) measure $P$.
}
  \onslide<3->{
 \item It can be shown that every random variable $X$ can be
 written as the difference of two nonnegative random variables.
}
  \onslide<4->{
 \item It can also be shown that every nonnegative random variable
 is the limit of a monotone sequence of simple random variables.
 Therefore, integration can be extended from $SRV(\Omega)$ to $RV(\Omega)$
}
\end{itemize}
\end{frame}

\begin{frame}
  \onslide<1->{
	\frametitle{Statistics} 
	\framesubtitle{Expectation}
}
\begin{itemize}
  \onslide<2->{
  \item Note that we did not need the use of a probability distribution nor density in defining the expectation this way.
}
  \onslide<3->{
  \item We only needed the probability function.
}
  \onslide<4->{
  \item Note also that the integration we defined above is abstract in the  sense that it works just as well for general measures, not just probability measures.
}
\end{itemize}
\end{frame}

\begin{frame}
  \onslide<1->{
	\frametitle{Statistics} 
	\framesubtitle{Expectation}
}
\begin{itemize}
  \onslide<2->{
 \item Expectation is a linear functional on the space $RV$ of random variables over $\Omega$:
 \[\expect{aX+bY}=a\expect{X}+b\expect{Y}\]
}
  \onslide<3->{
 \item Expectation is monotonic:
 \[X\leq Y a.s.\Rightarrow \expect{X}\leq \expect{Y}\]
}
  \onslide<4->{
 \item We can compute $\expect{X}$ by means of an ordinary
 Riemann-Stieltjes integral:
}
  \onslide<5->{
\begin{equation}\label{expect-RS}
  \expect{X}= \infint xdF(x)
\end{equation}
}
\end{itemize}
\end{frame}

\begin{frame}
  \onslide<1->{
	\frametitle{Statistics} 
	\framesubtitle{Expectation}
}
\begin{itemize}
  \onslide<2->{
 \item If $X$ also has a p.d.f, then equation \eqref{expect-RS}
 becomes:
\begin{equation}\label{expect-pdf}
  \expect{X} = \infint xdF(x) = \infint xF'(x)dx = \infint xf(x)dx
\end{equation}
}
  \onslide<3->{
 \item The expectation of $\exp(juX)$ for each $u\in \Real$ is
 called the \textit{\textbf{characteristic function}} of $X$:
\begin{equation}\label{charfunc}
  \phi_X(u)=\expect{\exp(juX)} = \infint \exp(jux)dF(x)
\end{equation}
where $j=\sqrt{-1}$.
}
\end{itemize}
\end{frame}
\begin{frame}
  \onslide<1->{
	\frametitle{Statistics} 
	\framesubtitle{2nd order moment:  Variance}
}
\begin{itemize}
  \onslide<2->{
 \item The variance of a random variable is defined as the mean square deviation from the mean:
\[\sigma^2 = \expect{(X-\mu)^2}\]
}
\end{itemize}
\end{frame}
\begin{frame}
  \onslide<1->{
	\frametitle{Statistics} 
	\framesubtitle{Higher Order Moments}
}
\begin{itemize}
  \onslide<2->{
 \item The value of the integral
 \[m_i=\infint x^nf(x)dx\]
 is called the $n^{th}$ \textit{\textbf{moment}} of the distribution.  If instead we
 subtract out the mean $\mu=\expect{X}$ prior to integrating, then we obtain the $n^{th}$
 \textit{\textbf{central moment}}:
 \[\mu_i=\infint (x-\mu)^nf(x)dx \]
}
  \onslide<3->{
 \item We note that $\mu=\mu_1$, the first order moment, or mean
}
\end{itemize}
\end{frame}
\begin{frame}
  \onslide<1->{
	\frametitle{Statistics} 
	\framesubtitle{Higher Order Moments}
}
\begin{itemize}
  \onslide<2->{
 \item The characteristic function $\phi_X$ is similar to the \textit{\textbf{moment generating
 function}} defined as follows:
\[M_X(t)=\expect{\exp(tX)}\]
}
  \onslide<3->{
 \item If $X$ has a density function $f$, then
 \[M_X(t)=\infint \exp(tx)f(x)dx\]
}
\end{itemize}
\end{frame}

\begin{frame}
  \onslide<1->{
	\frametitle{Statistics} 
	\framesubtitle{Higher Order Moments}
}
\begin{itemize}
  \onslide<2->{
 \item Note that
 \[\mu=\expect{X}=\frac{d}{dx}M_X(0)\]
  and, in general,
  \[\expect{X^n}=\frac{d^n}{dx^n}M_X(0).\]
}
  \onslide<3->{
 \item The Gaussian distribution has only two fundamental moments $\mu=\expect{X}$ and
 $\sigma^2=\expect{(X-\mu)^2}$.  All higher-order moments can be
 expressed in terms of these two.
}
  \onslide<4->{
 \item Thus, the Gaussian distribution is said to be a second-order distribution.
}
\end{itemize}
\end{frame}
\begin{frame}
  \onslide<1->{
	\frametitle{Statistics} 
	\framesubtitle{Covariance and Correlation Between Random Variables}
}
\begin{itemize}
  \onslide<2->{
\item  Let $X$ and $Y$ be two different random variables, 
i.e., functions from the sample space $\Omega$ into the reals $\Real$. 
}
  \onslide<3->{
\item Let $\sigsq{X}=\variance{X}$ be the variance of $X$ and $\sigsq{Y}=\variance{Y}$ be the variance of $Y$, 
and let $\mu_X,\mu_Y$ be their respective means.   
}
  \onslide<4->{
\item The \textbf{covariance} between $X$ and $Y$ is the mean of their product, 
after removing their means: $\covdef{X}{Y}$ 
}
  \onslide<5->{
\item Their correlation is their covariance normalized by the product of their standard deviations:
$\corrdef{X}{Y}$
}
\end{itemize}
\end{frame}


